% 摘要：仅陈述已完成并有运行记录支持的结果；问题二至四给出设计而非结果。
研究算力约束下提升大语言模型能力的资源配置建模。针对数据质量难以统一度量、指标间存在相互冲突、领域配比受单纯形约束、算力预算多维分配四类问题，本文以交叉熵损失核心测度指标，结合问题要求逐步建立“质量评价—冲突消解—配比建模—广义标度律—算力优化—前沿预测”的递进式分析体系。

问题一将原始数据进行预处理后结合参考文献选择了22个指标，采用 CRITIC评价模型客观赋权获得每个指标的权重、再使用加权距离 TOPSIS 获得样本级质量分，覆盖去重后 $261{,}086$ 条记录并聚合（聚合方式是什么？）至领域级与语料级。结果表明：抽样集与全量语料均分分别为 $60.95$ 与 $54.23$，降幅的约 $99.8\%$ 由领域构成变化解释；对比了六种评分方案间秩相关不低于 $0.967$，排序结论稳健。针对单一总分无法区分“优缺点相抵”型样本的盲区，本文以五维相对分位阈值定义候选冲突（抽样集冲突率 $29.52\%$，留出集 $29.47\%$），并推导出有限补偿消解模型的闭式解，补偿参数有明确上界含义；扩展集主要结论与抽样集一致。配比建模在 ilr 坐标下建立同尺度响应面，并通过重复设计的配对分析证明只含配比的跨尺度绝对损失预测结构不可辨识，据此将外推结论严格限定为相对排序。
\keywords{数据质量评价\quad CRITIC--TOPSIS\quad 质量冲突消解\quad 单纯形约束\quad 标度律\quad 算力资源配置}
